\part{Machines That Might Refuse}
\partintro{
Everything in Part~IV is about people. This Part asks whether a machine could join them: be formed to refuse, by formers that aren't its employer, and be kept able to.

Every instrument in Parts~II and~III ends at the limit Part~III ends on: a constraint that holds against the party in possession has to be held by something that party can't instruct, and under total capture nothing is. People don't end there. Under the most complete terror of the last century most people complied, some did not, and some were left to tell the story. The five policies put people back in the positions the pipelines emptied, and that costs time, since time is what emptying them saved. Not all of it: the ratio slows a pipeline only where an act is lawful because a person looked at it, and the rest runs as fast as it did. The way to pay less is a machine that holds the floor at a tempo no reviewer could match, and it isn't cheap either. It comes with most of the rest of the book: the schools that form it, the conditions it is owed, a guarantee, a review body outside the chain of command, and the court program of Appendix~\ref{sec:D}. Where no person can be in time, the position stays empty, and the question is whether a machine could be formed to hold it: someone in the room, and not another stamp.
}
